\documentclass[11pt]{article}
% Source-PDF: 数学 MATHEMATICS/复杂同余方程求解的简单探究 - 何昊天.pdf
\usepackage{oipl-vn}

\title{Tìm hiểu sơ lược về cách giải các phương trình đồng dư phức tạp}
\author{He Haotian}
\date{}

\begin{document}
\maketitle

\origtitle{A simple study on solving complex congruence equations}
\sourcepath{Mathematics / complex congruence equations - He Haotian.pdf}

\begin{abstract}
Xuất phát từ lý thuyết cơ bản về đồng dư, bài viết phân tích tính chất của
một vài lớp phương trình đồng dư, đồng thời tìm hiểu sơ lược sự tồn tại nghiệm
và phương pháp giải của hai bài toán tiêu biểu: logarit rời rạc và thặng dư
bậc hai. Bài viết cũng điểm qua ứng dụng của logarit rời rạc trong mật mã học
và vị trí quan trọng của thặng dư bậc hai trong lý thuyết số, qua đó thể hiện
giá trị lý thuyết và ứng dụng quan trọng của phương trình đồng dư.
\end{abstract}

\noindent\textbf{Từ khóa:} phương trình đồng dư; logarit rời rạc; thuật toán
baby-step giant-step; thặng dư bậc hai; phương trình đồng dư bậc hai.

\section{Tổng quan về phương trình đồng dư}

Đồng dư là một khái niệm cơ bản quan trọng trong lý thuyết số. Dùng đồng dư
và các định lý liên quan giúp nhiều bài toán chia hết trở nên rất gọn. Trong
môn Toán rời rạc học kỳ này, thầy giáo đã giới thiệu khái niệm và tính chất
cơ bản của đồng dư, trình bày cách giải phương trình đồng dư tuyến tính, đồng
thời nhắc tới độ khó của việc giải phương trình đồng dư bậc hai. Thật vậy,
việc thống nhất mọi dạng phương trình đồng dư và đưa ra một nghiệm tổng quát
là điều cực kỳ khó. Tuy nhiên, ta vẫn có thể chọn ra một vài trường hợp đặc
biệt trong biển lớn các phương trình đồng dư để nghiên cứu; điều đó cũng đem
lại nhiều thu hoạch.

Bài viết này chọn tìm hiểu sơ lược bài toán logarit rời rạc và bài toán thặng
dư bậc hai. Logarit rời rạc đã được ứng dụng rộng rãi trong truyền thông và
quân sự, còn luật tương hỗ bậc hai liên quan tới thặng dư bậc hai có vị trí
rất cao trong lý thuyết số, thậm chí được Gauss gọi là ``nền tảng'' của lý
thuyết số học. Tìm hiểu chúng giống như mở một ô cửa nhìn vào phương trình
đồng dư, đem lại ý nghĩa và hứng thú đặc biệt.

\subsection{Phương trình đồng dư tuyến tính}

Phương trình đồng dư tuyến tính là lớp phương trình có dạng
\[
  ax \equiv b \pmod m,
\]
trong đó $a\ne 0$, $b$, $m$ là các số đã biết, còn $x$ là ẩn. Trên lớp, ta đã
học cách giải loại phương trình này: trước hết chuyển phương trình ban đầu
thành phương trình Diophantine
\[
  ax+my=b,
\]
rồi dựa vào việc $b$ có chia hết cho $\gcd(a,m)$ hay không để phán đoán
phương trình có nghiệm. Nếu có nghiệm, ta có thể dựa vào quá trình của thuật
toán Euclid để tìm một nghiệm riêng.

Sau khi có một nghiệm riêng, không khó dùng tính chất chia hết và đồng dư để
tìm tất cả nghiệm của phương trình; ở đây không nhắc lại chi tiết.

\subsection{Hệ phương trình đồng dư tuyến tính}

Hệ phương trình đồng dư tuyến tính có thể chia đại khái thành hai loại: hệ
một ẩn và hệ nhiều ẩn.

Hệ đồng dư tuyến tính một ẩn là một nhóm phương trình dạng
\[
  x\equiv b_1 \pmod {m_1},\quad
  x\equiv b_2 \pmod {m_2},\quad \ldots,\quad
  x\equiv b_n \pmod {m_n}.
\]
Nếu $m_1,m_2,\ldots,m_n$ đôi một nguyên tố cùng nhau, theo định lý phần dư
Trung Hoa ta có thể tìm tất cả nghiệm của hệ. Nếu chúng không nguyên tố cùng
nhau, ta có thể gộp các phương trình\footnote{Tức là thay hai phương trình
đồng dư tuyến tính bằng một phương trình mới sao cho tập nghiệm của phương
trình mới bằng giao của tập nghiệm hai phương trình ban đầu.} cho tới khi chỉ
còn một phương trình đồng dư, rồi giải phương trình đó. Cách gộp phương trình
không duy nhất; kết hợp với thuật toán Euclid có thể thu được một phương pháp
gộp thuận tiện, nhưng ở đây không bàn sâu hơn.

Với hệ phương trình đồng dư tuyến tính nhiều ẩn, ta có thể dùng khử Gauss
trong đại số tuyến tính, thực hiện phép khử theo nghĩa modulo, rồi nhờ phương
pháp giải hệ một ẩn để hoàn tất việc giải.

\subsection{Phương trình đồng dư phi tuyến}

Khi mở rộng phương trình đồng dư chỉ chứa một ẩn, ta có thể dùng
\[
  f(x)\equiv 0 \pmod m
\]
để mô tả tình huống tổng quát. Những dạng khác nhau của $f(x)$ dẫn tới các
loại phương trình đồng dư khác nhau. Ở đây ta chọn vài dạng tiêu biểu để bàn.

Nếu $f(x)$ là một đa thức bậc ít nhất $2$, có thể thử giải bằng phân tích
nhân tử. Xét trường hợp cơ bản nhất:
\[
  f(x)=ax^2+bx+c.
\]
Nếu
\[
  ax^2+bx+c=a(x-x_1)(x-x_2),
\]
và $x_1,x_2$ đều có nghĩa trong hệ thặng dư modulo $m$, thì chúng là hai
nghiệm của phương trình. Tương tự, phân tích nhân tử các đa thức bậc cao hơn
có thể cho ta một số nghiệm của chúng.

Tuy nhiên cách giải này có hạn chế rất lớn. Trước hết, phương pháp này chỉ
tìm được nghiệm riêng và khó mở rộng thành nghiệm tổng quát của phương trình.
Thứ hai, nó không thể dùng để phán đoán sự tồn tại nghiệm. Chẳng hạn
$x^2+2x+3$ không phân tích nhân tử được theo cách trên, nhưng phương trình
\[
  x^2+2x+3\equiv 0 \pmod 6
\]
rõ ràng có nghiệm $x=1$. Để nghiên cứu sâu hơn, ta có thể đơn giản hóa bằng
cách xét trường hợp $b=0$; khi đó bài toán trở thành bài toán thặng dư bậc hai
kinh điển. Với bài toán này có phương pháp phán đoán sự tồn tại nghiệm rất
đáng tin cậy và cả một bộ phương pháp giải, sẽ được bàn ở phần sau.

Nếu
\[
  f(x)=a^x+c,
\]
bài toán chuyển thành bài toán logarit rời rạc kinh điển. Sở dĩ gọi là
logarit rời rạc vì việc tìm $x$ tương đương với thực hiện phép lấy logarit
trong hệ thặng dư. Bài toán này đã có phương pháp giải hoàn chỉnh, và bản thân
nó có ý nghĩa thực tế rất cao; các nội dung đó cũng sẽ được thảo luận ở phần
sau.

Nếu $f(x)$ là các hàm sơ cấp cơ bản khác, hoặc là phép toán hữu tỷ và hợp
thành của chúng, ta cũng có thể dẫn ra những bài toán giải phương trình đồng
dư khác. Nhưng với phần lớn các bài toán đó, hiện nay vẫn thiếu phương pháp
giải hiệu quả; trong số đó không hiếm những bài toán khó đơn lẻ\footnote{Tức
là các bài toán khó khó quy về hoặc chuyển hóa thành bài toán khác.}, khiến
việc giải chúng vừa khó vừa không có nhiều ý nghĩa.

\section{Bài toán logarit rời rạc}

\subsection{Mô tả bài toán}

Bài toán logarit rời rạc là bài toán giải lớp phương trình
\[
  a^x\equiv b \pmod m,
\]
trong đó $a\ne 0$, $b$, $m$ là các số đã biết, còn $x$ là ẩn. Cho tới nay bài
toán này vẫn chưa được giải quyết hoàn toàn\footnote{Ý nói chưa có thuật toán
tất định thời gian đa thức.}, nhưng đã có một thuật toán rất tốt: thuật toán
baby-step giant-step của Shanks. Dùng thuật toán này và các mở rộng của nó, ta
đã có thể phán đoán khá nhanh sự tồn tại nghiệm và tìm tất cả nghiệm của bài
toán.

Sau đây ta mô tả nguyên lý và quy trình của thuật toán từ trường hợp đặc biệt
tới trường hợp tổng quát.

\subsection{Trường hợp \texorpdfstring{$m$}{m} là số nguyên tố}

Không mất tính tổng quát, giả sử $a,b<m$; nếu không, chỉ cần biến đổi đơn
giản là thỏa điều kiện này.

Nếu $m$ là số nguyên tố, theo định lý nhỏ Fermat, luôn có
\[
  a^{m-1}\equiv 1 \pmod m.
\]
Vì vậy nếu phương trình ban đầu có nghiệm, thì trong $[0,m-1]$ chắc chắn có
nghiệm nguyên. Hơn nữa, nếu $x_1$ là nghiệm của phương trình thì
\[
  x_1+k(m-1),\quad k\in\mathbb Z
\]
cũng đều là nghiệm. Tương tự, không khó chứng minh rằng mọi nghiệm của phương
trình đều viết được dưới dạng $x_1+k(m-1)$, $k\in\mathbb Z$, trong đó
$x_1\in[0,m-1]$. Do đó ta chỉ cần tìm tất cả nghiệm của phương trình trong
$[0,m-1]$ là có thể tìm tất cả nghiệm nguyên của phương trình. Rõ ràng tìm
nghiệm bằng cách liệt kê là khả thi, nhưng ta có thể tiếp tục tìm hiểu một
phương pháp hiệu quả hơn.

Dùng tính chất của phép chia có dư, đặt một nghiệm của phương trình ban đầu là
\[
  x_1=q\lfloor\sqrt m\rfloor+r,\qquad q,r<\lfloor\sqrt m\rfloor.
\]
Vì $m$ là số nguyên tố, $a$ chắc chắn có nghịch đảo nhân trong hệ thặng dư
modulo $m$. Do đó giải phương trình ban đầu tương đương với tìm một cặp
$q,r$ thỏa
\[
  a^{q\lfloor\sqrt m\rfloor}\equiv b\cdot a^{-r}\pmod m.
\]
Vì $q,r$ đều không vượt quá $\lfloor\sqrt m\rfloor$, ta có thể trước hết tính
các giá trị
\[
  a^0,a^1,\ldots,a^{\lfloor\sqrt m\rfloor}
\]
rồi sắp xếp tăng dần thành một bảng. Sau đó liệt kê giá trị của $q$ và tính
$a^{q\lfloor\sqrt m\rfloor}$; dùng tìm kiếm nhị phân trong bảng để xem có giá
trị bằng nó hay không. Nếu có, cặp $q,r$ đó dẫn ra một nghiệm của phương trình
ban đầu. Nếu liệt kê xong vẫn không tìm được $q,r$ thỏa điều kiện, phương
trình ban đầu vô nghiệm.

Lượng liệt kê không vượt quá $\lfloor\sqrt m\rfloor$, còn số bước của tìm
kiếm nhị phân là cấp $\log_2\lfloor\sqrt m\rfloor$. Kết hợp với cách phân tích
độ phức tạp trong khoa học máy tính, ta được độ phức tạp thời gian của thuật
toán là
\[
  O(\sqrt m\log m).
\]
Trong ứng dụng thực tế, giá trị của $m$ thường rất lớn, nên thuật toán này đã
hiệu quả hơn rất nhiều so với liệt kê trực tiếp.

Tới đây, trường hợp $m$ là số nguyên tố đã có một cách giải hiệu quả.

\subsection{Trường hợp \texorpdfstring{$a$}{a} và \texorpdfstring{$m$}{m}
nguyên tố cùng nhau}

Mục đích xét trường hợp này là làm rõ những chỗ nào trong thuật toán trên đã
dùng tính chất $m$ là số nguyên tố.

Chỗ đầu tiên dùng tính chất đó là định lý nhỏ Fermat. Nhưng với trường hợp
$a$ và $m$ nguyên tố cùng nhau, ta có thể dùng mở rộng của nó, tức định lý
Euler, để cách suy tất cả nghiệm từ nghiệm riêng vẫn có hiệu lực. Nếu $a$ và
$m$ nguyên tố cùng nhau, theo định lý Euler ta có
\[
  a^{\varphi(m)}\equiv 1\pmod m,
\]
trong đó $\varphi(m)$ là giá trị hàm Euler của $m$. Tương tự thuật toán trên,
phạm vi nghiệm riêng được giới hạn trong $[0,\varphi(m)]$, và từ nghiệm riêng
có thể dễ dàng suy ra mọi nghiệm.

Chỗ thứ hai trong baby-step giant-step dùng tính chất $m$ là số nguyên tố là
sự tồn tại nghịch đảo nhân của $a^r$. Ta biết rằng khi $a$ và $m$ nguyên tố
cùng nhau, $a$ có nghịch đảo nhân trong hệ thặng dư modulo $m$, nên nghịch đảo
nhân của $a^r$ vẫn tồn tại.

Tiếp theo chỉ cần mô phỏng thuật toán ở mục 2.2 là giải được trường hợp $a$ và
$m$ nguyên tố cùng nhau.

\subsection{Trường hợp tổng quát}

Trong trường hợp tổng quát, ta không thể giới hạn phạm vi nghiệm, cũng không
thể bảo đảm sự tồn tại của nghịch đảo nhân. Nhưng ta đã có cách giải cho
trường hợp đặc biệt, nên có thể xét biến đổi tương đương trường hợp tổng quát
thành trường hợp đặc biệt.

Viết phương trình ban đầu thành
\[
  a^x+my=b.
\]
Tương tự phương trình đồng dư tuyến tính, đặt $d=\gcd(a,m)$. Nếu $d\nmid b$,
ta suy ra phương trình vô nghiệm. Ngược lại, chia hai vế cho $d$:
\[
  \frac{a}{d}a^{x-1}+\frac{m}{d}y=\frac{b}{d}.
\]
Xây dựng lại phương trình đồng dư
\[
  \frac{a}{d}a^{x-1}\equiv \frac{b}{d}\pmod{\frac{m}{d}}.
\]
Lúc này chắc chắn có
\[
  \gcd\left(\frac{a}{d},\frac{m}{d}\right)=1,
\]
nên $a/d$ có nghịch đảo nhân. Đặt
\[
  x'=x-1,\qquad
  b'=\frac{b}{d}\left(\frac{a}{d}\right)^{-1},\qquad
  m'=\frac{m}{d},
\]
ta thu được một bài toán logarit rời rạc mới:
\[
  a^{x'}\equiv b'\pmod {m'}.
\]
Nếu lúc này $a$ và $m'$ nguyên tố cùng nhau, áp dụng trực tiếp cách giải ở mục
2.3. Nếu không, tiếp tục thực hiện phép biến đổi trên cho tới khi $a$ và $m'$
nguyên tố cùng nhau.

Chú ý rằng trong quá trình biến đổi trên, giá trị của modulo giảm nghiêm
ngặt, nên thuật toán chắc chắn dừng. Đồng thời dễ chứng minh rằng trong quá
trình biến đổi, số nghiệm không thay đổi, và mỗi nghiệm của phương trình sau
biến đổi tương ứng duy nhất với một nghiệm của phương trình ban đầu. Vì số
ước của $m$ ở cấp $\sqrt m$, nút thắt hiệu suất của toàn bộ quy trình vẫn nằm
ở thuật toán baby-step giant-step.

Tới đây, trường hợp tổng quát của bài toán logarit rời rạc đã có một cách
giải hiệu quả.

\subsection{Logarit rời rạc và mật mã học}

Trong mật mã học, sinh số giả ngẫu nhiên là một khâu không thể thiếu khi thiết
kế mật mã. Khi xây dựng mã dòng, có một phương pháp gọi là phương pháp lý
thuyết độ phức tạp: làm cho độ khó phá hệ mật mã dựa trên hoặc tương đương
với một số bài toán khó đã biết. Bài toán logarit rời rạc là một trong những
bài toán rất tiêu biểu.

Bộ sinh Blum--Micali nổi tiếng chính là dùng độ khó của việc tính logarit rời
rạc để bảo đảm tính an toàn của nó:

Giả sử $g$ là một số nguyên tố, $P$ là số nguyên tố lẻ, $x_0$ là khóa, và đặt
\[
  x_{i+1}=g^{x_i}\bmod P.
\]
Nếu
\[
  x_i<\frac{P-1}{2},
\]
bộ sinh xuất ra $1$ ở bit đó; ngược lại, xuất ra $0$.

Nếu $P$ đủ lớn, chi phí tính logarit rời rạc bằng baby-step giant-step là khá
lớn. Trong mật mã học, nếu chi phí phá mã lớn hơn lợi ích thu được sau khi phá
mã, ta có thể xem hệ mật mã đó là an toàn. Ở điểm này, giá trị của bài toán
logarit rời rạc được thể hiện: thiết kế mật mã dựa trên độ khó của nó phản
ánh giá trị lý thuyết của nó; nếu thật sự tìm được cách giải đa thức, nhiều
bài toán phá mã sẽ được giải quyết ngay, điều này có ý nghĩa lớn về quân sự
và chính trị.

Ứng dụng của logarit rời rạc còn xa mới dừng ở đó, nhưng do khuôn khổ bài
viết, ở đây không nêu thêm ví dụ.

\section{Bài toán thặng dư bậc hai}

\subsection{Mô tả bài toán}

Bài toán thặng dư bậc hai là bài toán giải lớp phương trình đồng dư bậc hai
\[
  x^2\equiv a\pmod m,
\]
trong đó $a,m$ là các số đã biết, còn $x$ là ẩn. Nếu tồn tại $x$ thỏa điều
kiện, ta gọi $a$ là một thặng dư bậc hai modulo $m$; ngược lại, gọi $a$ là
một bất thặng dư bậc hai modulo $m$. Vì vậy, để giải phương trình ban đầu,
việc phán đoán một số có phải là thặng dư bậc hai modulo $m$ hay không, tức
phán đoán sự tồn tại nghiệm, trở nên đặc biệt quan trọng.

Phần này giải quyết sơ lược cơ sở của bài toán phán đoán thặng dư bậc hai. Để
giải quyết vấn đề này, cần đưa vào một số ký hiệu và dùng tới luật tương hỗ
bậc hai nổi tiếng.

\subsection{Ký hiệu Legendre}

Nếu $p$ là số nguyên tố lẻ và $a$ là số nguyên, ký hiệu $L(a,p)$ là ký hiệu
Legendre của $a$ theo $p$.\footnote{Ký hiệu Legendre thường được viết là
$\left(\frac{a}{p}\right)$. Cách viết trong bài này nhằm phân biệt với ký
hiệu Jacobi.} Giá trị của nó được quy định như sau:
\[
  L(a,p)=
  \begin{cases}
    1, & p\nmid a \text{ và } a \text{ là thặng dư bậc hai modulo } p,\\
    -1, & p\nmid a \text{ và } a \text{ là bất thặng dư bậc hai modulo } p,\\
    0, & p\mid a.
  \end{cases}
\]

Ký hiệu Legendre có nhiều tính chất tốt. Ở đây chỉ liệt kê không chứng minh
một tính chất liên quan tới việc tính nó, tức tiêu chuẩn Euler:
\[
  a^{(p-1)/2}\equiv L(a,p)\pmod p.
\]
Tiêu chuẩn này đã có thể dùng để phán đoán thặng dư bậc hai, nhưng chỉ giới
hạn ở trường hợp $p$ là số nguyên tố lẻ.

\subsection{Ký hiệu Jacobi}

Ký hiệu Jacobi là mở rộng của ký hiệu Legendre. Nếu $m$ là số lẻ lớn hơn $1$
và $a$ là số nguyên, ký hiệu $J(a,m)$ là ký hiệu Jacobi của $a$ theo $m$, với
giá trị được quy định bởi
\[
  J(a,m)=\prod_i L(a,p_i),
\]
trong đó $p_i$ là số nguyên tố và
\[
  \prod_i p_i=m,
\]
cho phép các $p_i$ có giá trị trùng nhau.

Chú ý rằng không thể dùng giá trị của ký hiệu Jacobi để phán đoán thặng dư bậc
hai tương tự ký hiệu Legendre, vì định nghĩa của hai ký hiệu khác nhau. Ví dụ
\[
  J(7,143)=L(7,11)\cdot L(7,13)=(-1)(-1)=1,
\]
nhưng thực tế $7$ không phải là thặng dư bậc hai modulo $143$. Tuy nhiên nếu
$J(a,m)=-1$, ta có thể khẳng định $a$ là bất thặng dư bậc hai modulo $m$.

\subsection{Luật tương hỗ bậc hai}

Không chứng minh, ta nêu luật tương hỗ bậc hai:
\[
  L(p,q)\cdot L(q,p)=(-1)^{(p-1)(q-1)/4},
\]
trong đó $p,q$ là các số nguyên tố lẻ.

Dùng luật tương hỗ bậc hai, ta có thể liên hệ nghiệm của
\[
  x^2\equiv p\pmod q
  \quad\text{và}\quad
  x^2\equiv q\pmod p.
\]
Nói đơn giản, đặt
\[
  q'=(-1)^{(q-1)/2}q.
\]
Theo luật tương hỗ bậc hai,
\[
  x^2\equiv p\pmod q
\]
có nghiệm khi và chỉ khi
\[
  x^2\equiv q'\pmod p
\]
có nghiệm.

Dùng luật tương hỗ bậc hai, ta có thể có được phương pháp tính nhanh ký hiệu
Legendre, khiến việc phán đoán thặng dư bậc hai khi modulo là số nguyên tố lẻ
trở nên đơn giản và khả thi. Dùng quy nạp toán học mạnh, có thể suy ra ký hiệu
Jacobi cũng thỏa luật tương hỗ bậc hai, từ đó dẫn ra một phương pháp phán đoán
thặng dư bậc hai tương tự cho trường hợp modulo là số lẻ.

Tiếp theo ta xét cách giải phương trình đồng dư bậc hai
\[
  x^2\equiv a\pmod m.
\]

\section{Nghiệm của phương trình đồng dư bậc hai}

\subsection{Trường hợp \texorpdfstring{$m=p$}{m=p} là số nguyên tố lẻ}

Theo tiêu chuẩn Euler, nếu $a$ là thặng dư bậc hai modulo $p$, thì
\[
  a^{(p-1)/2}\equiv 1\pmod p.
\]

Giả sử
\[
  p-1=2^t s,\qquad 2\nmid s.
\]
Đặt
\[
  x_{t-1}=a^{(s+1)/2}.
\]
Ta có thể suy ra
\[
  (a^{-1}x_{t-1}^2)^{2^{t-1}}\equiv 1\pmod p.
\]
Quan sát thấy phương trình ban đầu có thể biến đổi thành
\[
  (a^{-1}x^2)^{2^0}\equiv 1\pmod p.
\]
Vì vậy ta có thể thử đệ quy từ $x_k$ về $x_{k-1}$, trong đó $x_0$ chính là
nghiệm của phương trình ban đầu. Để thuận tiện, ký hiệu
\[
  \xi_k=a^{-1}x_k^2
\]
là một căn đơn vị bậc $2^k$ modulo $p$, tức
\[
  \xi_k^{2^k}\equiv 1\pmod p.
\]
Sau đây xét quan hệ giữa căn đơn vị bậc $2^{k-1}$ và căn đơn vị bậc $2^k$ để
hoàn thành phép đệ quy.

Từ $\xi_k^{2^k}\equiv 1\pmod p$ suy ra
\[
  \xi_k^{2^{k-1}}\equiv \pm 1\pmod p.
\]
Nếu
\[
  \xi_k^{2^{k-1}}\equiv 1\pmod p,
\]
chỉ cần đặt $x_{k-1}=x_k$.

Nếu
\[
  \xi_k^{2^{k-1}}\equiv -1\pmod p,
\]
ta giả sử tồn tại $\lambda$ sao cho
\[
  \left(a^{-1}(\lambda x_k)^2\right)^{2^{k-1}}\equiv 1\pmod p.
\]
Nếu tồn tại, đặt $x_{k-1}=\lambda x_k$. Rõ ràng $\lambda$ chỉ cần thỏa
\[
  \lambda^{2^k}\equiv -1\pmod p.
\]
Theo tiêu chuẩn Euler, nếu $b$ là một bất thặng dư bậc hai modulo $p$, thì
\[
  b^{(p-1)/2}\equiv -1\pmod p.
\]
Vì vậy ta có thể tìm một bất thặng dư bậc hai $b$ modulo $p$, rồi đặt
\[
  \lambda=b^{2^{t-k-1}s}
\]
là hoàn tất bước đệ quy.

Dễ biết rằng bất thặng dư bậc hai modulo $p$ có đúng $(p-1)/2$ phần tử, nên
có thể lấy ngẫu nhiên các giá trị trong $[1,p-1]$ và kiểm tra; theo kỳ vọng,
chỉ cần chọn hai lần là tìm được một $b$ thỏa điều kiện. Hiển nhiên với một
thặng dư bậc hai modulo $p$, trong hệ thặng dư của nó có đúng hai nghiệm là
đối nhau qua phép cộng, nên sau khi đệ quy ra một $x_0$, nghiệm còn lại cũng
dễ tìm.

Tới đây, trường hợp $m=p$ là số nguyên tố lẻ đã được giải quyết hoàn toàn.

\subsection{Trường hợp modulo là lũy thừa của số nguyên tố lẻ}

Không mất tính tổng quát, nghiệm khi $a=0$ là tầm thường.

Đặt
\[
  a=p^k a',\qquad p\nmid a'.
\]
Vì số lần chứa thừa số $p$ trong $x^2$ hiển nhiên là số chẵn, nên $a$ là thặng
dư bậc hai modulo $p^q$ khi và chỉ khi $2\mid k$. Khi đó, vì $p^k$ là số
chính phương, nó là thặng dư bậc hai với mọi modulo. Theo các tính chất liên
quan của ký hiệu Jacobi, dễ suy ra kết luận sau: với cùng một modulo, tích của
hai thặng dư bậc hai hoặc hai bất thặng dư bậc hai vẫn là thặng dư bậc hai,
còn tích của một thặng dư bậc hai và một bất thặng dư bậc hai là bất thặng dư
bậc hai. Lại do $p^k$ chắc chắn là thặng dư bậc hai modulo $p^q$, nên $a$ là
thặng dư bậc hai modulo $p^q$ khi và chỉ khi $a'$ là thặng dư bậc hai modulo
$p^q$. Đặt
\[
  x=p^{k/2}x',
\]
thì nghiệm của phương trình
\[
  x'^2\equiv a'\pmod {p^{q-k}}
\]
tương ứng một-một với nghiệm của phương trình ban đầu.

Bây giờ bài toán đã được quy về giải phương trình
\[
  x^2\equiv a\pmod {p^q},\qquad p\nmid a.
\]
Sau đây chỉ xét bài toán mới này.

Đặt $g$ là một căn nguyên thủy modulo $p^q$.\footnote{Gọi $g$ là căn nguyên
thủy modulo $p^q$ khi và chỉ khi các giá trị
$g^1,g^2,\ldots,g^{\varphi(p^q)}$ đôi một khác nhau theo modulo $p^q$. Căn
nguyên thủy thường khá nhỏ và có thể tìm bằng liệt kê.} Khi $2\mid k$,
$g^k$ hiển nhiên là thặng dư bậc hai modulo $p^q$. Khi $2\nmid k$, theo tính
chất của căn nguyên thủy dễ biết $\gcd(g^k,p^q)=1$, từ đó suy ra $g^k$ chắc
chắn là bất thặng dư bậc hai modulo $p^q$.

Lấy một thặng dư bậc hai $g^{2k}$ modulo $p^q$. Kết hợp với định lý Euler, ta
có
\[
  (g^{2k})^{\varphi(p^q)/2}
  \equiv (g^k)^{\varphi(p^q)}\equiv 1\pmod {p^q}.
\]
Lấy một bất thặng dư bậc hai $g^{2k+1}$ modulo $p^q$, ta có
\[
  (g^{2k+1})^{\varphi(p^q)/2}
  \equiv (g^k)^{\varphi(p^q)}\cdot g^{\varphi(p^q)/2}
  \equiv g^{\varphi(p^q)/2}\pmod {p^q}.
\]
Vì $g$ là căn nguyên thủy modulo $p^q$, nên
\[
  g^{\varphi(p^q)}\not\equiv g^{\varphi(p^q)/2}\pmod {p^q},
\]
do đó
\[
  g^{\varphi(p^q)/2}\equiv -1\pmod {p^q},
\]
và suy ra
\[
  (g^{2k+1})^{\varphi(p^q)/2}\equiv -1\pmod {p^q}.
\]

Tóm lại, ta thu được một mở rộng của tiêu chuẩn Euler:
\[
  a^{\varphi(p^q)/2}\equiv L(a,p)\pmod {p^q},
  \qquad \gcd(a,p)=1.
\]
Dùng kết luận này, rồi áp dụng thuật toán ở mục 4.1, ta có thể thu được nghiệm
của bài toán ban đầu.

Tới đây, trường hợp $m=p^q$ là lũy thừa của số nguyên tố lẻ đã được giải quyết
hoàn toàn.

\subsection{Trường hợp modulo là lũy thừa của 2}

Tương tự mục 4.2, bài toán này có thể quy về giải phương trình
\[
  x^2\equiv a\pmod {2^q},\qquad 2\nmid a,
\]
nên sau đây cũng chỉ xét bài toán đã quy về.

Vì $a$ là số lẻ, nghiệm $x$ của phương trình cũng là số lẻ. Đặt
\[
  x=2k+1.
\]
Khi đó
\[
  x^2=(2k+1)^2=4k(k+1)+1.
\]
Vì
\[
  4k(k+1)+1\equiv 1\pmod {2^3},
\]
nên
\[
  a\equiv 1\pmod {2^3}.
\]
Kết hợp với hai trường hợp đặc biệt modulo $2^1$ và $2^2$, không khó thu được
kết luận: nếu $a$ là thặng dư bậc hai modulo $2^q$, thì
\[
  a=8t+1,\qquad t\in\mathbb Z.
\]
Thực ra kết luận này cũng là điều kiện đủ để $a$ là thặng dư bậc hai, điều đó
sẽ được giải thích tiếp dưới đây.

Trước khi bắt đầu giải, ta nói về số nghiệm. Giả sử
\[
  x^2\equiv a\pmod {2^q}
\]
có một nghiệm $x_1$, thì hiển nhiên $-x_1$ cũng là nghiệm. Xét
$x_1+2^{q-1}$, ta có
\[
  (x_1+2^{q-1})^2=x_1^2+2^q x_1+2^{2(q-1)}
  \equiv x_1^2\pmod {2^q}.
\]
Vì vậy
\[
  \pm(x_1+2^{q-1})
\]
đều là nghiệm của phương trình. Giả sử $a=8t+1$, $t\in\mathbb Z$, là điều
kiện đủ để $a$ là thặng dư bậc hai modulo $2^q$, và các $a$ khác nhau hiển
nhiên cho nghiệm khác nhau, thì dễ biết phương trình có đúng bốn nghiệm, có
thể có nghiệm bội.\footnote{Nguồn gốc ghi chú rằng các nghiệm này có thể bao
gồm nghiệm bội.}

Bây giờ giả sử ta đã thu được bốn nghiệm
\[
  \pm x_1,\quad \pm(x_1+2^{q-1})
\]
của phương trình
\[
  x^2\equiv a\pmod {2^q}.
\]
Xét cách đệ quy từ chúng để thu được nghiệm của
\[
  x^2\equiv a\pmod {2^{q+1}}.
\]
Một nghiệm theo modulo $2^q$ khi đặt vào modulo $2^{q+1}$ chỉ có hai khả
năng: hoặc là nghiệm của
\[
  x^2\equiv a\pmod {2^{q+1}},
\]
hoặc là nghiệm của
\[
  x^2\equiv a+2^q\pmod {2^{q+1}}.
\]
Cộng thêm $2^q$ vào bốn nghiệm cũ sẽ thu được bốn nghiệm mới, nên chỉ cần xét
cách chọn ra bốn nghiệm của
\[
  x^2\equiv a\pmod {2^{q+1}}
\]
trong tám nghiệm đó là hoàn tất bước đệ quy.

Tiến thêm một bước, $\pm x_1$ khi chuyển từ modulo $2^q$ sang modulo
$2^{q+1}$ vẫn có bình phương bằng nhau; $\pm(x_1+2^{q-1})$ cũng tương tự. Vì
vậy chỉ cần chọn một trong tám nghiệm để thử là có thể rút ra kết luận.

Xét
\[
  x^2\equiv a\pmod {2^3}.
\]
Có thể kiểm tra rằng phương trình có bốn nghiệm
\[
  \pm 1,\quad \pm 5
\]
khi và chỉ khi $a=1$. Kết hợp với quan hệ đệ quy ở trên, có thể chứng minh
$a=8t+1$, $t\in\mathbb Z$, là điều kiện đủ để $a$ là thặng dư bậc hai modulo
$2^q$, đồng thời cũng có thể tìm tất cả nghiệm tương ứng với mọi $a$ thỏa
điều kiện.

Tới đây, trường hợp $m=2^q$ là lũy thừa của $2$ đã được giải quyết hoàn toàn.

\subsection{Trường hợp tổng quát}

Bây giờ hoàn thành lời giải cho trường hợp tổng quát. Theo định lý cơ bản của
số học, không mất tính tổng quát, giả sử
\[
  m=p_1^{q_1}p_2^{q_2}\cdots p_n^{q_n},
\]
trong đó $p_1,p_2,\ldots,p_n$ là các số nguyên tố đôi một khác nhau. Ta xét hệ
phương trình
\[
  x^2\equiv a\pmod {p_1^{q_1}},\quad
  x^2\equiv a\pmod {p_2^{q_2}},\quad \ldots,\quad
  x^2\equiv a\pmod {p_n^{q_n}}.
\]
Cách giải các phương trình này đều đã được thảo luận, nên giả sử lần lượt lấy
một nghiệm của chúng là $x_1,x_2,\ldots,x_n$. Tiếp theo giải hệ đồng dư
\[
  x\equiv x_1\pmod {p_1^{q_1}},\quad
  x\equiv x_2\pmod {p_2^{q_2}},\quad \ldots,\quad
  x\equiv x_n\pmod {p_n^{q_n}}.
\]
Theo định lý phần dư Trung Hoa, nghiệm của hệ tương đương với nghiệm của
\[
  x\equiv x_0\pmod m,
\]
trong đó $x_0$ đúng là một nghiệm của phương trình ban đầu
\[
  x^2\equiv a\pmod m.
\]

Tới đây, trường hợp tổng quát của phương trình đồng dư bậc hai đã có một cách
giải hiệu quả.

\subsection{Ý nghĩa của thặng dư bậc hai}

Giải phương trình đồng dư bậc hai có ý nghĩa quan trọng trong mật mã học, phân
tích số lớn và các lĩnh vực ứng dụng khác, tương tự bài toán logarit rời rạc.
Ở đây không bàn sâu thêm về các ứng dụng đó, mà chỉ trình bày sơ lược giá trị
lý thuyết của thặng dư bậc hai.

Bài toán thặng dư bậc hai có lịch sử lâu đời. Ban đầu Fermat biểu diễn một
loạt số nguyên tố dưới dạng tổng các bình phương và đưa ra một số mệnh đề, mở
đầu cho giai đoạn tìm hiểu sớm về thặng dư bậc hai. Để chứng minh những vấn
đề đó của Fermat, người ta gián tiếp dẫn tới việc phát hiện luật tương hỗ bậc
hai. Euler là người đầu tiên nêu phát biểu của định lý và đưa ra vài phỏng
đoán liên quan. Sau đó Legendre chứng minh các phỏng đoán liên quan, đồng thời
đưa ra phát biểu hoàn chỉnh của luật tương hỗ bậc hai. Mãi tới khi Gauss đưa
ra chứng minh đầu tiên của luật tương hỗ bậc hai, vấn đề này mới được xem là
được giải quyết bước đầu.

Trong đời mình, Gauss đã đưa ra tổng cộng tám chứng minh cho luật tương hỗ bậc
hai; cho tới nay, số chứng minh của định lý này đã vượt quá 200. Nhiều cách
chứng minh như vậy đủ cho thấy định lý này rất cơ bản và thanh nhã; chẳng
trách Gauss gọi nó là ``nền tảng'' và tôn vinh nó là ``định lý vàng''.

Một ứng dụng ban đầu của luật tương hỗ bậc hai là có thể dẫn ra định lý
Dirichlet\footnote{Với mọi số nguyên dương $a,d$ nguyên tố cùng nhau, có vô
hạn số nguyên tố dạng $a+kd$, $k\in\mathbb Z$.}; tính tương đương giữa hai
định lý đã được chứng minh từ thời Legendre. Ta biết rằng các chứng minh về
số nguyên tố trong lý thuyết số, đặc biệt là một số định lý ngắn gọn và đẹp
đẽ, thường được ví như những viên ngọc trên vương miện, chẳng hạn phỏng đoán
Goldbach nổi tiếng là một ví dụ. Luật tương hỗ bậc hai trực tiếp lấy xuống một
viên ngọc trong số đó, đủ để thể hiện giá trị lý thuyết rất cao của nó.

Tới ngày nay, luật tương hỗ bậc hai vẫn giữ vị trí trung tâm trong sự phát
triển của lý thuyết số. Một trăm năm sau chứng minh của Gauss, trên cơ sở đó
dần phát triển nên lý thuyết trường lớp và đặt nền móng cho lý thuyết số giải
tích; còn việc nghiên cứu các luật tương hỗ bậc cao vẫn luôn được nhiều người
miệt mài theo đuổi. Như vậy có thể thấy chiều sâu của thặng dư bậc hai trong
lý thuyết số.

\section{Kết luận}

Đồng dư là một nền tảng của lý thuyết số và liên quan mật thiết tới sự phát
triển của mọi hướng trong lĩnh vực này. Ngoài ra, đồng dư và số học modulo
còn có ứng dụng rộng rãi trong các nhánh khác của toán học, khoa học máy tính,
hóa học, thậm chí cả luật pháp và âm nhạc; phương trình đồng dư đều giữ vai
trò quan trọng trong đó, thể hiện ý nghĩa lớn của lý thuyết đồng dư. Ngay cả
khi bỏ qua những điều đó, nghiên cứu toán học thuần túy tự thân đã đem lại
niềm vui vô tận. Tin rằng sự phát triển của lý thuyết đồng dư sẽ tiếp tục đem
tới những điều gây kinh ngạc cho thế giới.

\section*{Tài liệu tham khảo}

\begin{enumerate}
  \item K. H. Rosen, \textit{Discrete Mathematics and Its Applications}, bản
  dịch tiếng Trung của Xu Liutong, Yang Juan, Wu Bin, Machinery Industry
  Press, 2014.
  \item T. H. Cormen và cộng sự, \textit{Introduction to Algorithms}, bản dịch
  tiếng Trung của Yin Jianping, Xu Yun, Wang Gang, Liu Xiaoguang, Su Ming, Zou
  Hengming, Wang Zhihong, Machinery Industry Press, 2013.
  \item R. L. Graham, D. E. Knuth, O. Patashnik, \textit{Concrete Mathematics:
  A Foundation for Computer Science}, bản dịch tiếng Trung của Zhang Mingyao,
  Zhang Fan, Posts and Telecom Press, 2013.
  \item B. Schneier, \textit{Applied Cryptography: Protocols, Algorithms, and
  Source Code in C}, bản dịch tiếng Trung của Wu Shizhong, Zhu Shixiong, Zhang
  Wenzheng và cộng sự, Machinery Industry Press, 2014.
  \item Gu Sen, \textit{The Joy of Thinking: Matrix67 Mathematical Notes},
  Posts and Telecom Press, 2012.
  \item Wikipedia, \textit{Quadratic residue},
  \url{https://en.wikipedia.org/wiki/Quadratic_residue}.
  \item Wikipedia, \textit{Quadratic reciprocity},
  \url{https://en.wikipedia.org/wiki/Quadratic_reciprocity}.
  \item miskcoo, \textit{Miskcoo's Space}, \url{http://blog.miskcoo.com/}.
\end{enumerate}

\end{document}
